\documentclass[a4paper]{article}

% PACKAGES %%%%%%%%%%%%%%%%%%%%%%%%%%%%%%%%%%%%%%%%%%%%%%%%%%%%%%%%%%%%%%%%%%%%%
\usepackage[utf8]{inputenc}
\usepackage[english]{babel}
\usepackage{url}
\usepackage{hyperref}
\usepackage{charter}
\usepackage{float}
\usepackage{fancyhdr}
\usepackage{booktabs}
\usepackage{tabularx}
\usepackage{multirow}
\usepackage{enumitem}
\usepackage{xcolor}
\usepackage{graphicx}
\usepackage{siunitx}
\usepackage{comment}
\usepackage[hmarginratio=1:1, top=32mm, margin=0.8in, headsep=20mm]{geometry}

\newcommand{\TODO}{\textbf{\textcolor{red}{TODO}}}

\lhead{\includegraphics[scale=0.4]{CETC.png}}
\rhead{\includegraphics[scale=0.025]{Cornell-University-Logo.png}}

\pagestyle{fancy}
\fancyfoot{}
\fancyfoot[C]{\thepage}

\begin{document}

\thispagestyle{fancy}

\begin{center}
    \textbf{\large Orbit Odyssey -- Computer Vision Autonomous Challenge} \\[0.3em]
    {\Large Quotation}

    \vspace{1em}

    David Alejandro Fuquen Fl\'orez \\[0.3em]
    {\small \today}
\end{center}

\vspace{1em}

%===============================================================================
\section*{Summary}
%===============================================================================

This quotation covers the remaining development work for the Autonomous Computer Vision portion of the Orbit Odyssey Challenge, originally proposed by Rushil Sambangi (Cornell University). The scope includes three deliverables: a complete competition manual, adaptation of the graphical training interface, and refinement of the XRP robot control code.

The project will be delivered in three sequential phases.

%===============================================================================
\section*{Hourly Rate}
%===============================================================================

\begin{center}
\begin{tabular}{ll}
\toprule
\textbf{Hourly rate} & \$28.00 USD \\
\textbf{Estimated total hours} & 45 hours \\
\textbf{Estimated total cost} & \$1,260.00 USD \\
\bottomrule
\end{tabular}
\end{center}

%===============================================================================
\section*{Phase 1: Competition Manual}
%===============================================================================

\noindent\textbf{Deadline:} 26/05/2026 (to be shown at 10 AM HKT meeting) \hfill \textbf{Estimated hours:} 10 \hfill \textbf{Cost:} \$280.00

\vspace{0.5em}

The manual will serve as the single source of truth for the competition. It will be inspired by and expand upon the original \href{https://docs.google.com/document/d/1N_mCX1EM3w5nDMbHrX3vcQlSqa2yJL38nCU0xxMPxmY/edit?tab=t.0}{Orbit Odyssey Manual}, with a strong emphasis on preventing cheating and closing loopholes.

\subsection*{Scope of Work}

\begin{enumerate}[leftmargin=*]
    \item \textbf{Playing Field Definition} (3 hrs)
    \begin{itemize}
        \item Field dimensions, boundaries, and markings
        \item Robot starting zone placement rules
        \item Object quantity, placement templates, and clearance constraints
    \end{itemize}

    \item \textbf{Scoring System and Scorecards} (3 hrs)
    \begin{itemize}
        \item Complete point system (field performance, mAP bonus, model-size multiplier)
        \item Penalty definitions and edge-case rulings
        \item Tiebreaker hierarchy
        \item Printable scorecard template for referees
        \item Integration of cryptographic receipt verification into scoring workflow (see Phase 2 section for details)
    \end{itemize}

    \item \textbf{Student Guide} (1.5 hrs)
    \begin{itemize}
        \item Step-by-step instructions for using the training GUI
        \item Explanation of hyperparameter choices and their trade-offs
        \item How to submit the trained model and receipt for deployment
    \end{itemize}

    \item \textbf{Moderator Guide} (1.5 hrs)
    \begin{itemize}
        \item Field setup checklist (per-match randomization)
        \item Receipt verification procedure (decryption and validation) (see Phase 2 section for details)
    \end{itemize}

    \item \textbf{Anti-Cheating Documentation} (1 hr)
    \begin{itemize}
        \item How the signed receipt system works (for judges)
        \item What constitutes a valid vs.\ tampered receipt
        \item Disqualification criteria
    \end{itemize}
\end{enumerate}

\subsection*{Hour Breakdown}

\begin{center}
\begin{tabular}{lcc}
\toprule
\textbf{Task} & \textbf{Hours (est.)} & \textbf{Cost} \\
\midrule
Playing field definition & 3.0 & \$84.00 \\
Scoring system \& scorecards & 3.0 & \$84.00 \\
Student guide & 1.5 & \$42.00 \\
Moderator guide & 1.5 & \$42.00 \\
Anti-cheating documentation & 1.0 & \$28.00 \\
\midrule
\textbf{Phase 1 Total} & \textbf{10.0} & \textbf{\$280.00} \\
\bottomrule
\end{tabular}
\end{center}

%===============================================================================
\section*{Phase 2: Graphical User Interface}
%===============================================================================

\noindent\textbf{Timeline:} June 9--12, 2026 \hfill \textbf{Estimated hours:} 22 \hfill \textbf{Cost:} \$616.00

\vspace{0.5em}

This phase adapts the existing training GUI (\texttt{gui\_train\_model.py}) to work with CETC's current codebase. It also implements a cryptographic anti-cheating system: a signed, encrypted receipt that records the exact hyperparameters and results from each training run. Only moderators possess the decryption key.

\subsection*{Scope of Work}

\begin{enumerate}[leftmargin=*]
    \item \textbf{CETC Code Integration} (9 hrs)
    \begin{itemize}
        \item Analyze CETC's existing training/data pipeline
        \item Adapt GUI to interface with CETC's code structure
        \item Ensure compatibility with CETC's dataset format and workflow
        \item Resolve any dependency or conflicts
    \end{itemize}

    \item \textbf{Signed Receipt System} (7 hrs)
    \begin{itemize}
        \item Implement Fernet symmetric encryption with a moderator-held secret key
        \item Receipt captures: hyperparameters, class names, final mAP per class, model file size, timestamp, and unique run ID
        \item Receipt exported as encrypted \texttt{.receipt} file alongside model weights
        \item Moderator-side decryption utility for verification
    \end{itemize}

    \item \textbf{mAP and Model-Size Extraction} (2 hrs)
    \begin{itemize}
        \item Programmatically extract mAP from training results CSV
        \item Compute and record model file size (bytes)
        \item Feed both values into the receipt and display in GUI summary
    \end{itemize}

    \item \textbf{UI Adjustments and Testing} (4 hrs)
    \begin{itemize}
        \item Update dataset selection flow for CETC's folder structure
        \item Add receipt export status indicator in GUI
        \item End-to-end testing: train $\rightarrow$ export $\rightarrow$ receipt generation $\rightarrow$ decryption verification
    \end{itemize}
\end{enumerate}

\subsection*{Hour Breakdown}

\begin{center}
\begin{tabular}{lcc}
\toprule
\textbf{Task} & \textbf{Hours (est.)} & \textbf{Cost} \\
\midrule
CETC code analysis \& integration & 9.0 & \$252.00 \\
Signed receipt system (Fernet) & 7.0 & \$196.00 \\
mAP / model-size extraction & 2.0 & \$56.00 \\
UI adjustments \& testing & 4.0 & \$112.00 \\
\midrule
\textbf{Phase 2 Total} & \textbf{22.0} & \textbf{\$616.00} \\
\bottomrule
\end{tabular}
\end{center}

%===============================================================================
\section*{Phase 3: XRP Robot Code}
%===============================================================================

\noindent\textbf{Timeline:} June 19--23, 2026 \hfill \textbf{Estimated hours:} 13 \hfill \textbf{Cost:} \$364.00

\vspace{0.5em}

Once the playing field and game rules are finalized (Phase 1), the XRP robot's MicroPython control code (\texttt{main.py}) will be adapted to operate within the defined arena constraints. Note: Given hardware constraints (I do not possess the camera or chip to be used for this challenge) this part must be done in partnership with CETC's engineers for proper testing.

\subsection*{Scope of Work}

\begin{enumerate}[leftmargin=*]
    \item \textbf{State Machine Adaptation} (4 hrs)
    \begin{itemize}
        \item Align search/center/contact/avoid behavior with defined field dimensions
        \item Parameterize hardcoded values (distances, speeds, delays)
        \item Handle edge cases
    \end{itemize}

    \item \textbf{Parameter Tuning and Testing} (5 hrs)
    \begin{itemize}
        \item Tune turn speeds, centering threshold, approach distances for physical field
        \item Test with actual objects (cans, wheels, rocket model)
        \item Validate rangefinder accuracy at competition distances
    \end{itemize}

    \item \textbf{Integration Testing} (4 hrs)
    \begin{itemize}
        \item Full pipeline test: GUI-trained model $\rightarrow$ \texttt{run\_model.py} $\rightarrow$ XRP
        \item Verify ACK flow control stability over full 120-second run
        \item Confirm scoring conditions are met
    \end{itemize}
\end{enumerate}

\subsection*{Hour Breakdown}

\begin{center}
\begin{tabular}{lcc}
\toprule
\textbf{Task} & \textbf{Hours (est.)} & \textbf{Cost} \\
\midrule
State machine adaptation & 4.0 & \$112.00 \\
Parameter tuning \& testing & 5.0 & \$140.00 \\
Integration testing (full pipeline) & 4.0 & \$112.00 \\
\midrule
\textbf{Phase 3 Total} & \textbf{13.0} & \textbf{\$364.00} \\
\bottomrule
\end{tabular}
\end{center}

%===============================================================================
\section*{Total Cost Summary}
%===============================================================================

\begin{center}
\begin{tabular}{lccc}
\toprule
\textbf{Phase} & \textbf{Hours} & \textbf{Cost} & \textbf{Deadline} \\
\midrule
1. Competition Manual & 10.0 & \$280.00 & May 26, 2026 \\
2. GUI Adaptation & 22.0 & \$616.00 & June 12, 2026 \\
3. XRP Robot Code & 13.0 & \$364.00 & June 23, 2026 \\
\midrule
\textbf{Total} & \textbf{45.0} & \textbf{\$1,260.00} & \\
\bottomrule
\end{tabular}
\end{center}

%===============================================================================
\section*{Milestones}
%===============================================================================

\begin{description}[leftmargin=!,labelwidth=5em]
    \item[May 26, 2026] Complete competition manual. Deliver: finalized PDF with playing field specs, scoring rules, student/moderator guides, anti-cheat documentation, and printable scorecards.

    \item[June 9--12, 2026] CETC's training code integration and receipt system. Deliver: working GUI that produces encrypted receipts alongside trained models. Moderator decryption utility functional.

    \item[June 19--23, 2026] XRP code finalization and full-pipeline integration testing. Deliver: updated \texttt{main.py}, tested end-to-end with physical hardware and competition objects.
\end{description}

%===============================================================================
\subsection*{Itinerary and Limitations}
%===============================================================================

\begin{itemize}
    \item \textbf{June 1--7, 2026:} Unavailable due to a work trip. No work will be performed during this period. This gap is accounted for in the timeline: Phase 1 concludes before June 1, and Phase 2 begins on June 9.
    \item \textbf{Hardware dependency (Phase 3):} I do not currently possess the XRP robot, webcam, or AIxBoard required for physical testing. Phase 3 must be conducted in partnership with CETC engineers who have access to the hardware. Delays in hardware access may affect the June 23 deadline.
\end{itemize}

%===============================================================================
\section*{Payment Terms}
%===============================================================================

Payment may be structured per-phase or upon project completion, as agreed with the team. Hours will be tracked and reported weekly.

\end{document}
